\chapter{The Information Plane in Practice}
\label{chap:info_plane_practice}

In previous chapters, we formalized the Information Bottleneck principle (Equation~\ref{eq:ib_lagrangian}) and explored the theoretical bounds of mutual information in deep neural networks (Theorem~\ref{thm:mi_bound}). To bridge the gap between theory and empirical observation, we will train a simple 3-layer multi-layer perceptron (MLP) on the MNIST dataset and track its layer-wise information trajectories. We begin by defining the network architecture and modifying its forward pass to return not only the final predictions but also the continuous activations $T_i$ at each hidden layer, providing direct access to these representations.

\begin{pycode}
import torch
import torch.nn as nn
from torchvision import datasets, transforms
from torch.utils.data import DataLoader

class MLP(nn.Module):
    def __init__(self):
        super().__init__()
        self.fc1 = nn.Linear(784, 256)
        self.act1 = nn.Tanh()
        self.fc2 = nn.Linear(256, 128)
        self.act2 = nn.Tanh()
        self.fc3 = nn.Linear(128, 10)
        
    def forward(self, x):
        x = x.view(-1, 784)
        t1 = self.act1(self.fc1(x))
        t2 = self.act2(self.fc2(t1))
        y = self.fc3(t2)
        return y, [t1, t2]

# Prepare data and model
transform = transforms.Compose([transforms.ToTensor(), transforms.Normalize((0.5,), (0.5,))])
dataset = datasets.MNIST(root='./data', train=True, download=True, transform=transform)
loader = DataLoader(dataset, batch_size=256, shuffle=True)
model = MLP()
optimizer = torch.optim.SGD(model.parameters(), lr=0.1)
criterion = nn.CrossEntropyLoss()
\end{pycode}

With our network defined, the next challenge is estimating the mutual information quantities $I(X; T_i)$ and $I(T_i; Y)$. Estimating $I(X; T_i)$ is famously tricky in deterministic neural networks since continuous variables mapped through deterministic functions theoretically possess infinite mutual information, unless bounded by noise or binning. A robust workaround, recalling our discussions in Section~\ref{sec:mi_estimation_challenges} and Chapter~\ref{ch:mi_estimators}, is to use the KSG estimator on the multi-dimensional vectors. We add a small amount of Gaussian noise to identical points (such as categorical labels) to make the distance calculations tractable, and then apply the k-nearest neighbors method on the unreduced representations.

\begin{pycode}
import numpy as np
from scipy.special import digamma
from scipy.spatial import cKDTree

def compute_mi_ksg(T, Y, k=5):
    """
    Estimates Mutual Information I(T; Y) using the KSG k-nearest neighbor method.
    Works directly on the unreduced multi-dimensional vectors.
    """
    T = T.detach().cpu().numpy()
    Y = Y.detach().cpu().numpy()
    
    if len(T.shape) == 1:
        T = T[:, None]
    if len(Y.shape) == 1:
        Y = Y[:, None]
        
    # Add small noise to handle identical points
    T = T + 1e-8 * np.random.randn(*T.shape)
    Y = Y + 1e-8 * np.random.randn(*Y.shape)
        
    n = T.shape[0]
    TY = np.concatenate([T, Y], axis=1)
    
    tree_ty, tree_t, tree_y = cKDTree(TY), cKDTree(T), cKDTree(Y)
    distances, _ = tree_ty.query(TY, k=k+1, p=np.inf)
    epsilons = distances[:, k]
    
    # Count points in marginal spaces within epsilon distances
    nt = np.array([len(tree_t.query_ball_point(T[i], epsilons[i], p=np.inf)) - 1 for i in range(n)])
    ny = np.array([len(tree_y.query_ball_point(Y[i], epsilons[i], p=np.inf)) - 1 for i in range(n)])
    
    # KSG Formula
    mi = digamma(k) - np.mean(digamma(nt + 1) + digamma(ny + 1)) + digamma(n)
    return max(0.0, float(mi))

def estimate_information_plane(model, loader):
    model.eval()
    all_x, all_y, all_t1, all_t2 = [], [], [], []
    for x, y in loader:
        with torch.no_grad():
            _, (t1, t2) = model(x)
            all_x.append(x.view(x.size(0), -1))
            all_y.append(y)
            all_t1.append(t1)
            all_t2.append(t2)
            
    X_full = torch.cat(all_x)
    Y_full = torch.cat(all_y)
    T1_full = torch.cat(all_t1)
    T2_full = torch.cat(all_t2)
    
    # MI computations
    ix_t1 = compute_mi_ksg(X_full, T1_full)
    ix_t2 = compute_mi_ksg(X_full, T2_full)
    
    it1_y = compute_mi_ksg(T1_full, Y_full)
    it2_y = compute_mi_ksg(T2_full, Y_full)
    
    return [ix_t1, ix_t2], [it1_y, it2_y]
\end{pycode}

We can now execute the training loop over several epochs, systematically evaluating the mutual information coordinates at the end of each pass. As training progresses, the Information Bottleneck theory posits that the network undergoes two distinct phases: an initial rapid "fitting" phase where $I(X; T_i)$ increases as the network memorizes the input structure, followed by a slower "compression" phase where $I(X; T_i)$ decreases as the network sheds task-irrelevant nuisance variables. Concurrently, the predictive information regarding the targets, $I(T_i; Y)$, should rise and eventually plateau as cross-entropy loss is minimized.

\begin{pycode}
epochs = 15
trajectories_x = {1: [], 2: []}
trajectories_y = {1: [], 2: []}

for epoch in range(epochs):
    model.train()
    for batch_x, batch_y in loader:
        optimizer.zero_grad()
        out, _ = model(batch_x)
        loss = criterion(out, batch_y)
        loss.backward()
        optimizer.step()
        
    # Evaluate Information Plane coordinates
    ix_t, it_y = estimate_information_plane(model, loader)
    
    trajectories_x[1].append(ix_t[0])
    trajectories_x[2].append(ix_t[1])
    trajectories_y[1].append(it_y[0])
    trajectories_y[2].append(it_y[1])
    
    print(f"Epoch {epoch+1} completed: Layer 1 I(X;T)={ix_t[0]:.2f}")
\end{pycode}

Finally, we visualize the collected trajectories on the Information Plane. The x-axis represents the compression term $I(X; T_i)$—how much information the representation retains about the raw input—and the y-axis represents the prediction term $I(T_i; Y)$—how much information the representation contains about the desired output. Each curve corresponds to a specific hidden layer, mapping its journey epoch by epoch. When properly estimated, this plot typically reveals deeper layers positioned further to the left and higher up, visually confirming that they achieve greater compression and have more successfully distilled the task-relevant features.

\begin{pycode}
import matplotlib.pyplot as plt

plt.figure(figsize=(8, 6))

# Plot Layer 1 and 2 trajectories
plt.plot(trajectories_x[1], trajectories_y[1], marker='o', 
         linestyle='-', color='blue', label='Layer 1')
plt.plot(trajectories_x[2], trajectories_y[2], marker='s', 
         linestyle='-', color='orange', label='Layer 2')

plt.title('Information Plane Trajectories on MNIST')
plt.xlabel('I(X; T) [bits]')
plt.ylabel('I(T; Y) [bits]')
plt.legend()
plt.grid(True)
plt.show()
\end{pycode}
